\documentclass{article}
\usepackage[T1]{fontenc}
\usepackage{lmodern}
\usepackage{graphicx}
\usepackage{amsmath,amssymb}
\usepackage[a4paper,margin=2cm]{geometry}
\usepackage[absolute,overlay]{textpos}
\setlength{\TPHorizModule}{1mm}
\setlength{\TPVertModule}{1mm}
\usepackage[indonesian]{babel}
\usepackage{hyperref}
\setlength{\emergencystretch}{2em}
% unit: Halaman 45

\begin{document}

% === Halaman 45 (llm) ===

\begin{itemize}
    \item Alice memilih bilangan acak $k$ yang harus terletak di dalam selang $[1, p-1]$
    \item Cipherteks dari $P_M$ adalah pasangan titik
    \begin{itemize}
        \item $P_C = [ ( kB), (P_M + kP_B) ]$
    \end{itemize}
    \item Untuk mendekripsi, Bob mula-mula menghitung hasil kali titik pertama dari $P_C$ (yaitu $kB$) dengan kunci privatnya, $b$
    \begin{itemize}
        \item $b \cdot (kB)$
    \end{itemize}
    \item Bob kemudian mengurangkan titik kedua dari $P_C$ dengan hasil kali di atas
    \begin{itemize}
        \item $(P_M + kP_B) - [b \cdot (kB)] = P_M + k \cdot (bB) - b \cdot (kB) = P_M$
    \end{itemize}
    \item Bob kemudian men-decode $P_M$ untuk memperoleh pesan M
\end{itemize}

\vspace{10mm}

\textcolor{red}{*) Sumber bahan: Debdeep Mukhopadhyay, Elliptic Curve Cryptography , Dept of Computer Sc and Engg IIT Madras}

\end{document}
